\section{The triangular module and the counterexample}
\label{sec:triangular}

We now combine the two twists and apply the transfer principle of
Proposition~\ref{tr:transfer}.  Retain $E,X$ and the finite
side-projective bimodule $Y$ from Proposition~\ref{b:bimodule}.
For brevity, write $LX=L\otimes_E X$ for an $E$-bimodule $L$, and set
\[
 W^a=\sHom_E(X,YX[a]).
\]
Thus $W^a=k$ for $a=-3,0$ and is zero otherwise.

\subsection{The fiber of the two evaluated maps}

For $i=1,2$, put $U_i=U_{H_i}$.  Choose the maps
$g_i:U_i\to Y$ supplied by Proposition~\ref{b:bimodule}.
Using the fixed identifications $U_iX\simeq X$, their evaluated stable
classes are nonzero elements of the same one-dimensional space
$W^0$.  Rescale them so that both classes equal a fixed $w\ne0$.
Choose a finite free $E$-bimodule $Q$ with a surjection
$\pi:Q\to Y$, and define $F$ by
\begin{equation}\label{t:fiber}
 0\longrightarrow F\longrightarrow U_1\oplus U_2\oplus Q
 \xrightarrow{(g_1,g_2,\pi)}Y\longrightarrow0.
\end{equation}
For example, the action gives such a surjection from $E^e\otimes_kY$.
Since $Y$ is projective on each side, this sequence splits separately
on the left and on the right.  In particular, $F$ is finite and
side-projective.

Let
\[
 \mathcal V^a=\sHom_E(X,FX[a]).
\]
Evaluating \eqref{t:fiber} on $X$ preserves exactness, and $QX$ is
$E$-projective.  The stable long exact sequence is therefore
\begin{equation}\label{t:fiber-hom}
 (\cH^{a-1})^2\longrightarrow W^{a-1}
 \longrightarrow\mathcal V^a
 \xrightarrow{p^a}(\cH^a)^2\longrightarrow W^a.
\end{equation}
In degree zero the last map is
$(c,d)\mapsto(c+d)w$.  It follows that $\mathcal V^0$ is a line with a
unique element $v_0$ projecting to $(1,1)$.

For a stable map $g:X\to X[a]$, let $F(g)$ denote its image under
the exact functor $F\otimes_E-$, using the shift comparisons of
Lemma~\ref{h:evaluation}.  Define
\begin{equation}\label{t:delta}
 \Delta^a:(\cH^a)^2\longrightarrow\mathcal V^a,
 \qquad (g,h)\longmapsto F(g)v_0+v_0[a]h.
\end{equation}

\begin{lemma}\label{t:comparison-maps}
The map $\Delta^0$ is surjective.  For every $a>0$ and every
$a\leq-3$, $\Delta^a$ is an isomorphism, and $\Delta^{-2}$ is
injective.
\end{lemma}

\begin{proof}
The projection $p^a$ in \eqref{t:fiber-hom} is an isomorphism for
$a>0$.  For $a\geq2$ this follows from the support of $W$; at $a=1$
it also uses the surjectivity of $(\cH^0)^2\to W^0$.
The same support argument applies for $a\leq-4$.
At $a=-3$, both $W^{-4}$ and $\cH^{-3}$ are zero, so
$\mathcal V^{-3}=0$ and $p^{-3}$ is again an isomorphism.
The remaining relevant low degrees can be read directly from
\eqref{t:fiber-hom}:
\begin{equation}\label{t:boundary-fiber}
 \mathcal V^{-2}=k,
 \qquad \mathcal V^0=kv_0,
 \qquad \mathcal V^1=0.
\end{equation}

Naturality of the bimodule projections $F\to U_i$ gives
\begin{equation}\label{t:projected-delta}
 p^a\Delta^a(g,h)=
 \bigl(U_1(g)+h,\ U_2(g)+h\bigr).
\end{equation}
Indeed each projection sends $v_0$ to the identity of $X$, and
evaluation of a bimodule map commutes with the specified shift
comparisons.  By Proposition~\ref{e:twists}, at $a=3m>0$
the map in \eqref{t:projected-delta} is the scalar matrix
\[
 \begin{pmatrix}H_1^{-m}&1\\H_2^{-m}&1\end{pmatrix}
 \otimes 1_{\mathcal P_m},
\]
and at $a=-3m-1\leq-4$ it is
\[
 \begin{pmatrix}H_1^{m+2}&1\\H_2^{m+2}&1\end{pmatrix}
 \otimes 1_{D\mathcal P_m}.
\]
In both cases $m\geq1$.  Algebraic independence of $H_1,H_2$
implies that each determinant is nonzero in $k$, for every such $m$.
In all other degrees in the asserted ranges, the source and target
are zero by Proposition~\ref{e:all-degrees} and the already proved
isomorphism $p^a$.

Finally, $\Delta^0(c,d)=(c+d)v_0$ is surjective.  Since
$\cH^{-2}=0$, the map $\Delta^{-2}:0\to\mathcal V^{-2}=k$ is
injective.  Its surjectivity is neither claimed nor needed.
\end{proof}

\subsection{A cone over the triangular algebra}

Form the finite-dimensional algebra
\begin{equation}\label{t:triangular-algebra}
 \Lambda=\begin{pmatrix}E&0\\F&E\end{pmatrix}.
\end{equation}
A left $\Lambda$-module is a triple
$(L,N,\alpha:F\otimes_E L\to N)$ of left $E$-modules and an
$E$-linear structure map, with action
\[
 \begin{pmatrix}a&0\\b&d\end{pmatrix}(x,y)
 =\bigl(ax,\alpha(b\otimes x)+dy\bigr).
\]
The column constructions are standard in the theory of totally acyclic
complexes over triangular matrix rings; see
\cite[Lemma~2.2 and Proposition~2.3]{EshraghiEtAl2014v1}.
Here $F$ is projective on each side, so the required tensor exactness
and orthogonality conditions hold. We retain the complete resolution
to compute its self-Hom groups. Define the two column functors
\[
 \mathsf D_0(L)=(L,FL,1),
 \qquad \mathsf D_1(L)=(0,L,0).
\]
They are exact and preserve finite projectives: exactness uses right
projectivity of $F$, and their values on $E$ are the two columns of
the regular module $\Lambda$.  The descriptions of module maps give
\begin{align}
 \Hom_\Lambda(\mathsf D_i(L),\mathsf D_i(L'))
   &\simeq\Hom_E(L,L')\quad(i=0,1),\label{t:diagonal-hom}\\
 \Hom_\Lambda(\mathsf D_0(L),\mathsf D_1(L'))&=0,
 \qquad
 \Hom_\Lambda(\mathsf D_1(L),\mathsf D_0(L'))
   \simeq\Hom_E(L,FL').\label{t:cross-hom}
\end{align}
In the first diagonal slot, the map on the lower component is $F(g)$
when the map on the upper component is $g$.

Let $p_f=f\otimes f$.  Use the embedding
$i_X:X\hookrightarrow Ep_f$, $i_X(1)=\zeta$, established in
Lemma~\ref{e:socle}.  Choose a complete resolution
$P$ of $X$ over $E$ whose first negative term is
$P_{-1}=Ep_f$, and whose degree-zero differential is
$d_0=i_X\epsilon$, where $\epsilon:P_0\to X$ is the augmentation.
Such a choice is possible because $E$ is symmetric.  The complex
$FP$ is exact with projective terms, and $c_0(FP)=FX$.

Choose an ordinary map $\widetilde v:X\to FX$ representing $v_0$.
Lemma~\ref{h:comparison} lifts it to a chain map $P\to FP$.
By \eqref{t:cross-hom} this is a chain map
\[
 V:\mathsf D_1(P)\longrightarrow\mathsf D_0(P).
\]
Set
\begin{equation}\label{t:cone}
 P_Z=\Cone(V),\qquad Z=c_0P_Z.
\end{equation}
Thus
\[
 (P_Z)_j=\mathsf D_0(P_j)\oplus\mathsf D_1(P_{j-1}),
\]
with cross differential $V_{j-1}$ from the second summand to the
first summand in degree $j-1$.

\begin{lemma}\label{t:finite-module}
The complex $P_Z$ is totally acyclic.  Its degree-zero cokernel $Z$
is finite-dimensional and nonprojective.
\end{lemma}

\begin{proof}
Both column complexes have projective terms and are exact.  Since
$\Lambda=\mathsf D_0(E)\oplus\mathsf D_1(E)$,
\eqref{t:diagonal-hom}--\eqref{t:cross-hom} identify their Hom
complexes into $\Lambda$ with sums of
\[
 \Hom_E(P,E),\qquad\Hom_E(P,F),\qquad0.
\]
These are exact: $P$ is totally acyclic and $F$ is finite projective
on the left.  Hence both column complexes are totally acyclic, and
so is their cone.

For an explicit finite description, $Z$ is the triple with upper
component $X$ and lower component
\begin{equation}\label{t:explicit-module}
 N_Z=(FX\oplus Ep_f)
 \big/\{(\widetilde v(x),i_X(x)):x\in X\},
\end{equation}
whose structure map $FX\to N_Z$ is induced by the first summand.
To verify this formula, first quotient the degree-zero lower
component by the image of the differential in $\mathsf D_0(P)$.
This gives $FX$ in its first summand.  The remaining relations arise
from $P_0$ in the shifted second column.  Its two components induce
$\widetilde v\epsilon$ and $i_X\epsilon$, respectively, yielding
\eqref{t:explicit-module}.

The upper-component functor sends the projective generators
$\mathsf D_0(E)$ and $\mathsf D_1(E)$ to $E$ and $0$, and therefore
sends every finite projective $\Lambda$-module to an $E$-projective.
But $X$ is nonprojective, since its degree-three self-extension group
is nonzero by Proposition~\ref{e:all-degrees}.  Thus $Z$ is
nonprojective.
\end{proof}

\subsection{The full endomorphism complex}

The complete resolution makes it possible to check both vanishing
ranges with the same endomorphism calculation.

\begin{proposition}\label{t:vanishing}
The module and resolution in \eqref{t:cone} satisfy
\[
 H^a\Hom_\Lambda(P_Z,Z)=0
 \qquad(a>0\ \text{or}\ a\leq-2).
\]
\end{proposition}

\begin{proof}
Write $A_0=\mathsf D_0(P)$ and $B_0=\mathsf D_1(P)$, so that the
underlying graded complex of $P_Z$ is $A_0\oplus B_0[1]$.
By \eqref{t:cross-hom}, the block
$\Hom_\Lambda^\bullet(A_0,B_0[1])$ is zero.  Projection onto the two
diagonal blocks therefore gives a short exact sequence of full Hom
complexes
\begin{equation}\label{t:end-sequence}
\begin{split}
 0\longrightarrow\mathcal C^\bullet
 \longrightarrow\End_\Lambda^\bullet(P_Z)
 \longrightarrow
 \End_\Lambda^\bullet(A_0)\oplus
 \End_\Lambda^\bullet(B_0[1])
 \longrightarrow0,\\[-2pt]
 \mathcal C^\bullet=\Hom_\Lambda^\bullet(B_0[1],A_0).
\end{split}
\end{equation}
The underlying degreewise Hom spaces are products as in
Lemma~\ref{h:comparison}; the block decomposition has only four
blocks.  In particular, it does not replace any product by a direct
sum.

The diagonal cohomology is $(\cH^a)^2$.  To fix the cross-block
index explicitly, in degree $a$ we have
\begin{align*}
 \mathcal C^a
 &=\prod_j\Hom_E(P_{j-1},FP_{j-a})\\
 &=\prod_i\Hom_E(P_i,FP_{i-(a-1)})
   =\Hom_E^{a-1}(P,FP).
\end{align*}
The differentials agree under this reindexing in characteristic two.
Lemma~\ref{h:comparison} consequently gives
$H^a\mathcal C^\bullet=\mathcal V^{a-1}$.

The connecting map in \eqref{t:end-sequence}, from diagonal
cohomology in degree $a$ to $H^{a+1}\mathcal C^\bullet=\mathcal V^a$,
is precisely $\Delta^a$.  Indeed, lift two diagonal cocycles to the
diagonal blocks of an endomorphism of the cone and take its
cochain differential.  The cross component is the sum of its two
compositions with $V$.  The first diagonal acts downstairs by $F$
of the corresponding map, while the second acts by precomposition.
Under the composition-compatible comparison in
Lemma~\ref{h:comparison}, this gives
$F(g)v_0+v_0[a]h$, as in \eqref{t:delta}.
Thus the long exact sequence gives
\begin{equation}\label{t:end-cohomology}
 0\longrightarrow\coker\Delta^{a-1}
 \longrightarrow H^a\End_\Lambda^\bullet(P_Z)
 \longrightarrow\ker\Delta^a\longrightarrow0.
\end{equation}

For $a>1$ or $a\leq-3$, both outer terms vanish by
Lemma~\ref{t:comparison-maps}.  At $a=1$, they vanish by surjectivity
of $\Delta^0$ and injectivity of $\Delta^1$.  At $a=-2$, they vanish
by the isomorphism $\Delta^{-3}$ and injectivity of $\Delta^{-2}$.
Finally, apply Lemma~\ref{h:comparison} to the totally acyclic
complex $P_Z$ to identify this cohomology with
$H^a\Hom_\Lambda(P_Z,Z)$.
\end{proof}

\begin{remark}\label{t:neighboring-degrees}
The boundary degrees can also be located completely.
Equation~\eqref{t:fiber-hom} gives
$\mathcal V^{-1}=(\cH^{-1})^2=k^2$.
By Proposition~\ref{e:twists}, the two rows of $\Delta^{-1}$ are
$(H_i^2,1)$, so this map is invertible.  Since $\Delta^{-2}:0\to k$
and $\ker\Delta^0=k$, \eqref{t:end-cohomology} yields
\[
 H^a\End_\Lambda^\bullet(P_Z)=
 \begin{cases}
 k,&a=-1,0,\\
 0,&a\notin\{-1,0\}.
 \end{cases}
\]
In particular the nonzero cokernel of $\Delta^{-2}$ contributes in
degree $-1$, while degree $-2$ vanishes as required by
Proposition~\ref{tr:transfer}.
\end{remark}

\subsection{Completion of the proof}

\begin{proof}[Proof of Theorem~\ref{thm:main}]
Fix the finite bimodule maps in \eqref{t:fiber}, the complete
resolution $P$, and the lifted map used in \eqref{t:cone}.
These choices exist over $k$ by the preceding propositions.
The algebra $\Lambda$ in \eqref{t:triangular-algebra} and its module
$Z$ in \eqref{t:explicit-module} are finite-dimensional.
By Lemma~\ref{t:finite-module}, $Z$ is nonprojective and has the
totally acyclic resolution $P_Z$.
Proposition~\ref{t:vanishing} supplies both hypotheses on its
complete self-Hom groups in Proposition~\ref{tr:transfer}.
Apply that proposition with $R=\Lambda$ to obtain
\[
 A=\Lambda\ltimes D\Lambda,
 \qquad M=A\otimes_\Lambda Z.
\]
It proves that $A$ is finite-dimensional and symmetric, that $M$
is finite-dimensional and nonprojective, and that
$\Ext_A^i(M,M)=0$ for every $i>0$.  These are all the assertions of
Theorem~\ref{thm:main}.
\end{proof}
